
\documentclass[11pt]{article}

\usepackage[margin=1in]{geometry}
\usepackage{amsmath,amssymb,amsthm,mathtools}
\usepackage{enumitem}
\usepackage{hyperref}
\usepackage[nameinlink,capitalize]{cleveref}

\hypersetup{colorlinks=true,linkcolor=blue,citecolor=blue,urlcolor=blue}

\newtheorem{theorem}{Theorem}
\newtheorem{proposition}{Proposition}
\newtheorem{lemma}{Lemma}
\newtheorem{corollary}{Corollary}
\newtheorem{definition}{Definition}
\newtheorem{remark}{Remark}

\newcommand{\Z}{\mathbb Z}
\newcommand{\Q}{\mathbb Q}
\newcommand{\F}{\mathbb F}
\newcommand{\E}{\mathbb E}
\newcommand{\G}{\Gamma}
\newcommand{\1}{\mathbf 1}
\newcommand{\ord}{\operatorname{ord}}
\newcommand{\lcm}{\operatorname{lcm}}
\newcommand{\TV}{\operatorname{TV}}
\newcommand{\KL}{D_{\operatorname{KL}}}

\title{A Subgroup Obstruction Calculus for Two-Generator Exponential Congruences}
\author{Jared Wilder}
\date{Draft for submission --- Oracle-polished version --- 2026-05-13}

\begin{document}
\maketitle

\begin{abstract}
We develop an elementary finite-group calculus for local congruence obstructions in two-generator exponential problems of the form
\[
m a^k b^\ell+1.
\]
If \(a,b,m\) are units modulo \(d\), local solvability is governed exactly by membership of \(-m^{-1}\) in the subgroup
\[
\Gamma_d(a,b)=\langle a,b\rangle\subset(\Z/d\Z)^*.
\]
This gives an exact local density, a complete two-point distribution law for the local obstruction random variable, all raw moments, exact information-theoretic distances from uniformity, a finite Fourier expansion through the character annihilator of \(\Gamma_d(a,b)\), exact centered moment laws, and CRT orthogonality for weighted prime-level sums.  The motivating example is the Erd\H{o}s--Graham Problem \#203 sequence \(m2^k3^\ell+1\), but all results in this paper are elementary finite algebra and probability.  No prime-distribution theorem and no resolution of Erd\H{o}s--Graham Problem \#203 is claimed.
\end{abstract}

\section{Introduction}

Let \(a,b\in\Q^*\), and consider the congruence
\[
m a^k b^\ell+1\equiv0\pmod d.
\]
When \(a,b,m\) are units modulo \(d\), this is equivalent to
\[
a^k b^\ell\equiv -m^{-1}\pmod d.
\]
Thus local solvability is not mysterious: it is exactly a subgroup-membership question in the finite abelian group \((\Z/d\Z)^*\).

For Erd\H{o}s--Graham Problem \#203, the bases are \(a=2\), \(b=3\), and \((m,6)=1\).  At a modulus \(d\) coprime to \(6m\), the local subgroup is
\[
\Gamma_d=\langle2,3\rangle\subset(\Z/d\Z)^*.
\]
The exact local density is
\[
\rho_m(d)=\frac{\1_{\Gamma_d}(-m^{-1})}{|\Gamma_d|}.
\]
This paper records the finite algebra behind that formula and its consequences.

The central point is that a normalized subgroup indicator carries an entire obstruction calculus: exact distribution, all moments, character expansion, centered variance, and CRT orthogonality.  These statements are elementary, but they create a clean interface between local algebra and later analytic or sieve questions.

\subsection*{Main theorem package}

The paper proves the following finite statements.

\begin{enumerate}[leftmargin=1.5em]
\item A finite abelian subgroup obstruction lemma: normalized subgroup indicators have exact mean, variance, and character expansion.
\item A complete two-point distribution law for the local obstruction random variable.
\item Exact raw moments and exact information-theoretic distances from uniformity.
\item A multiplication-fiber theorem for the map \(H_1\times H_2\to H_1H_2\).
\item An exact local density formula for \(m a^k b^\ell+1\) modulo \(d\).
\item Exact moment and weighted-variance identities over the parameter \(m\).
\end{enumerate}

No analytic number theory is required for any of these results.

\section{Finite abelian subgroup obstruction calculus}

\begin{definition}
Let \(G\) be a finite abelian group and \(\Gamma\le G\).  Define
\[
\Delta_\Gamma(c)=\frac{\1_\Gamma(c)}{|\Gamma|}-\frac1{|G|},\qquad c\in G.
\]
We call \(\Delta_\Gamma\) the centered subgroup obstruction attached to \(\Gamma\).
\end{definition}

\begin{theorem}[Finite abelian subgroup obstruction lemma]\label{T:finite}
For a finite abelian group \(G\) and subgroup \(\Gamma\le G\),
\[
\sum_{c\in G}\Delta_\Gamma(c)=0,
\]
and
\[
\sum_{c\in G}|\Delta_\Gamma(c)|^2=\frac1{|\Gamma|}-\frac1{|G|}.
\]
Moreover, if \(\widehat G\) is the character group and
\[
\Gamma^\perp=\{\chi\in\widehat G:\chi(\gamma)=1\ \forall \gamma\in\Gamma\},
\]
then
\[
\Delta_\Gamma(c)=
\frac1{|G|}
\sum_{\substack{\chi\in\Gamma^\perp\\ \chi\ne1}}\chi(c).
\]
\end{theorem}

\begin{proof}
The mean-zero identity follows from
\[
\sum_{c\in G}\frac{\1_\Gamma(c)}{|\Gamma|}=1
\qquad\text{and}\qquad
\sum_{c\in G}\frac1{|G|}=1.
\]
For the second moment,
\[
\sum_{c\in G}\left(\frac{\1_\Gamma(c)}{|\Gamma|}-\frac1{|G|}\right)^2
=
|\Gamma|\left(\frac1{|\Gamma|}-\frac1{|G|}\right)^2
+(|G|-|\Gamma|)\frac1{|G|^2},
\]
which simplifies to \(1/|\Gamma|-1/|G|\).

For the character expansion, finite Fourier inversion gives
\[
\frac{\1_\Gamma(c)}{|\Gamma|}
=
\frac1{|G|}\sum_{\chi\in\Gamma^\perp}\chi(c).
\]
Subtracting \(1/|G|\), the contribution of the trivial character, gives the claimed formula.
\end{proof}

\begin{definition}[Annihilator index and obstruction energy]
Define the annihilator index
\[
\kappa_\Gamma=\frac{|G|}{|\Gamma|}=|\Gamma^\perp|
\]
and the obstruction energy
\[
\mathcal E(\Gamma)=\frac1{|\Gamma|}-\frac1{|G|}
=
\frac{\kappa_\Gamma-1}{|G|}.
\]
\end{definition}

\begin{corollary}[Energy monotonicity]\label{C:energy-monotone}
If \(\Gamma_1\le\Gamma_2\le G\), then
\[
\mathcal E(\Gamma_2)\le \mathcal E(\Gamma_1).
\]
\end{corollary}

\begin{proof}
This is immediate from \(|\Gamma_1|\le|\Gamma_2|\).
\end{proof}

\section{Complete distribution law}

The variance identity in \cref{T:finite} is the second-order shadow of a complete elementary distribution law.

\begin{theorem}[Scaled Bernoulli law]\label{T:scaled-bernoulli}
Let \(G\) be a finite abelian group, \(\Gamma\le G\), and let \(c\) be uniformly distributed on \(G\).  Define
\[
X_\Gamma(c)=\frac{\1_\Gamma(c)}{|\Gamma|}.
\]
Then
\[
X_\Gamma=
\begin{cases}
|\Gamma|^{-1},&\text{with probability }|\Gamma|/|G|,\\
0,&\text{with probability }1-|\Gamma|/|G|.
\end{cases}
\]
Consequently, for every integer \(r\ge1\),
\[
\E X_\Gamma^r
=
\frac1{|G|\,|\Gamma|^{r-1}}.
\]
\end{theorem}

\begin{proof}
The random variable \(X_\Gamma\) equals \(|\Gamma|^{-1}\) exactly on the subgroup \(\Gamma\), which has probability \(|\Gamma|/|G|\), and equals \(0\) elsewhere.  The moment formula follows immediately.
\end{proof}

\begin{corollary}[Information-theoretic obstruction profile]\label{C:info}
Let
\[
U_\Gamma(c)=\frac{\1_\Gamma(c)}{|\Gamma|},
\qquad
U_G(c)=\frac1{|G|}.
\]
Then
\[
\|U_\Gamma-U_G\|_{\TV}
=
1-\frac{|\Gamma|}{|G|},
\]
\[
\chi^2(U_\Gamma\|U_G)
=
\frac{|G|}{|\Gamma|}-1,
\]
and
\[
\KL(U_\Gamma\|U_G)
=
\log\frac{|G|}{|\Gamma|}.
\]
\end{corollary}

\begin{proof}
For total variation,
\[
\frac12\sum_{c\in G}\left|U_\Gamma(c)-U_G(c)\right|
=
\frac12\left[
|\Gamma|\left(\frac1{|\Gamma|}-\frac1{|G|}\right)
+
(|G|-|\Gamma|)\frac1{|G|}
\right]
=
1-\frac{|\Gamma|}{|G|}.
\]
For chi-square divergence,
\[
\chi^2(U_\Gamma\|U_G)
=
\sum_{c\in G}\frac{(U_\Gamma(c)-U_G(c))^2}{U_G(c)}
=
|G|\left(\frac1{|\Gamma|}-\frac1{|G|}\right)
=
\frac{|G|}{|\Gamma|}-1.
\]
For KL divergence,
\[
\KL(U_\Gamma\|U_G)
=
\sum_{c\in\Gamma}\frac1{|\Gamma|}
\log\left(\frac{1/|\Gamma|}{1/|G|}\right)
=
\log\frac{|G|}{|\Gamma|}.
\]
\end{proof}

\begin{remark}
The obstruction strength is therefore measured exactly by the subgroup index \(|G|/|\Gamma|\).  Larger generated subgroups are closer to uniform in every metric above.
\end{remark}

\section{Multiplication fibers}

\begin{lemma}[Multiplication fiber theorem]\label{L:fiber}
Let \(G\) be an abelian group, and let \(H_1,H_2\le G\).  For \(c\in G\),
\[
|\{(x,y)\in H_1\times H_2:xy=c\}|
=
|H_1\cap H_2|\ \1_{H_1H_2}(c).
\]
\end{lemma}

\begin{proof}
The multiplication map
\[
\mu:H_1\times H_2\to H_1H_2,\qquad (x,y)\mapsto xy
\]
is surjective.  Its kernel is
\[
\{(u,u^{-1}):u\in H_1\cap H_2\},
\]
which has size \(|H_1\cap H_2|\).  Hence every element of \(H_1H_2\) has exactly \(|H_1\cap H_2|\) preimages, and elements outside \(H_1H_2\) have none.
\end{proof}

\section{Two-generator exponential congruences}

Let \(d\ge1\), and suppose \(a,b,m\in(\Z/d\Z)^*\).  Define
\[
G_d=(\Z/d\Z)^*,\qquad \Gamma_d(a,b)=\langle a,b\rangle\subset G_d.
\]

\begin{proposition}[Exact local density]\label{P:localdensity}
The proportion of exponent pairs in the finite orbit \(\langle a\rangle\times\langle b\rangle\) satisfying
\[
m a^k b^\ell+1\equiv0\pmod d
\]
is
\[
\rho_{m,a,b}(d)=
\frac{\1_{\Gamma_d(a,b)}(-m^{-1})}{|\Gamma_d(a,b)|}.
\]
\end{proposition}

\begin{proof}
The congruence is equivalent to \(a^k b^\ell\equiv -m^{-1}\pmod d\).  Apply \cref{L:fiber} to
\[
H_1=\langle a\rangle,\qquad H_2=\langle b\rangle.
\]
The number of pairs mapping to a target \(c\) is \(|H_1\cap H_2|\1_{\Gamma_d(a,b)}(c)\).  Dividing by \(|H_1||H_2|\) and using
\[
|\Gamma_d(a,b)|=\frac{|H_1||H_2|}{|H_1\cap H_2|}
\]
gives the formula.
\end{proof}

\section{Moments over \texorpdfstring{\(m\)}{m}}

For \(m\in G_d\), define
\[
\Delta_m(d;a,b)=
\frac{\1_{\Gamma_d(a,b)}(-m^{-1})}{|\Gamma_d(a,b)|}
-
\frac1{\varphi(d)}.
\]

\begin{theorem}[Exact moment laws]\label{T:moments}
For every \(d\) with \(a,b\in G_d\),
\[
\sum_{m\in G_d}\Delta_m(d;a,b)=0,
\]
and
\[
\sum_{m\in G_d}|\Delta_m(d;a,b)|^2
=
\frac1{|\Gamma_d(a,b)|}-\frac1{\varphi(d)}.
\]
\end{theorem}

\begin{proof}
The map \(m\mapsto -m^{-1}\) is a bijection of \(G_d\).  The result is therefore \cref{T:finite} with \(G=G_d\) and \(\Gamma=\Gamma_d(a,b)\).
\end{proof}

\section{Character expansion}

\begin{theorem}[Character-annihilator expansion]\label{T:char}
For \(G_d=(\Z/d\Z)^*\) and \(\Gamma_d(a,b)=\langle a,b\rangle\),
\[
\Delta_m(d;a,b)=
\frac1{\varphi(d)}
\sum_{\substack{\chi\bmod d\\ \chi\ne\chi_0\\ \chi(a)=\chi(b)=1}}
\chi(-m^{-1}).
\]
\end{theorem}

\begin{proof}
Characters \(\chi\bmod d\) are the character group of \(G_d\).  The annihilator of \(\Gamma_d(a,b)\) is exactly the set of characters satisfying \(\chi(a)=\chi(b)=1\).  Apply \cref{T:finite}.
\end{proof}

\section{CRT orthogonality and weighted variance}

\begin{lemma}[CRT orthogonality]\label{L:crt}
Let \(d=d_1d_2\) with \((d_1,d_2)=1\).  If \(F_i\) has mean zero on \((\Z/d_i\Z)^*\), then
\[
\E_{m\in(\Z/d\Z)^*}F_1(m\bmod d_1)F_2(m\bmod d_2)=0.
\]
\end{lemma}

\begin{proof}
By the Chinese remainder theorem,
\[
(\Z/d\Z)^*\simeq(\Z/d_1\Z)^*\times(\Z/d_2\Z)^*,
\]
and the expectation factors into the product of the two means.
\end{proof}

\begin{theorem}[Prime-level weighted variance]\label{T:weighted}
Let \(\mathcal Q\) be a finite set of primes not dividing \(ab\), and put \(P=\prod_{q\in\mathcal Q}q\).  Then
\[
\E_{m\in(\Z/P\Z)^*}
\left|
\sum_{q\in\mathcal Q}\lambda_q\Delta_m(q;a,b)
\right|^2
=
\sum_{q\in\mathcal Q}
|\lambda_q|^2
\frac1{q-1}
\left(
\frac1{|\Gamma_q(a,b)|}-\frac1{q-1}
\right).
\]
\end{theorem}

\begin{proof}
Expand the square.  The cross terms vanish by \cref{L:crt} and the mean-zero identity.  The diagonal terms are the normalized versions of \cref{T:moments} at prime modulus \(q\).
\end{proof}

\section{Specialization to \texorpdfstring{\(a=2,b=3\)}{a=2, b=3}}

For \(q\nmid6\), write
\[
\Gamma_q=\langle2,3\rangle\subset\F_q^*,
\qquad
a_q=\ord_q(2),\quad b_q=\ord_q(3).
\]
Since \(\F_q^*\) is cyclic,
\[
|\Gamma_q|=\lcm(a_q,b_q).
\]
For \(c\in\F_q^*\),
\[
|\{(x,y)\in\langle2\rangle\times\langle3\rangle:xy=c\}|
=
\gcd(a_q,b_q)\1_{\Gamma_q}(c).
\]
Thus for \((m,6q)=1\),
\[
\rho_m(q)=\frac{\1_{\Gamma_q}(-m^{-1})}{|\Gamma_q|},
\]
and
\[
\sum_{m\in\F_q^*}\left|
\rho_m(q)-\frac1{q-1}
\right|^2
=
\frac1{|\Gamma_q|}-\frac1{q-1}.
\]

\section{RMS obstruction law}

At prime modulus \(q\nmid ab\), \cref{T:moments} gives the normalized RMS size
\[
\mathrm{RMS}_m(\Delta_m(q;a,b))
=
\left[
\frac1{q-1}\left(
\frac1{|\Gamma_q(a,b)|}-\frac1{q-1}
\right)
\right]^{1/2}.
\]
Thus the local obstruction is governed exactly by the subgroup index
\[
[\F_q^*:\Gamma_q(a,b)].
\]
If \(|\Gamma_q(a,b)|\asymp q\), then the RMS size is \(O(q^{-1})\).  If the subgroup is smaller, the RMS size is amplified by
\[
\left(\frac{q-1}{|\Gamma_q(a,b)|}\right)^{1/2}.
\]

\section{Multi-generator extension}

The same formalism applies to \(t\)-generator exponential congruences.  Let
\[
a_1,\dots,a_t\in(\Z/d\Z)^*
\]
and put
\[
\Gamma_d(a_1,\dots,a_t)=\langle a_1,\dots,a_t\rangle\subset(\Z/d\Z)^*.
\]
For congruences
\[
m a_1^{n_1}\cdots a_t^{n_t}+1\equiv0\pmod d,
\]
the exact local density is
\[
\rho_{m,a_1,\dots,a_t}(d)
=
\frac{\1_{\Gamma_d(a_1,\dots,a_t)}(-m^{-1})}{|\Gamma_d(a_1,\dots,a_t)|},
\]
and all moment and character-expansion identities above remain valid with \(\Gamma_d(a,b)\) replaced by \(\Gamma_d(a_1,\dots,a_t)\).

\section{Concrete examples}

These examples are finite checks and are included to make the obstruction calculus tangible.

\subsection{No obstruction when \texorpdfstring{\(\Gamma_q=\F_q^*\)}{Gamma\_q = F\_q\^{}*}}

For \(q=5\),
\[
\F_5^*=\{1,2,3,4\},\qquad \langle2,3\rangle=\F_5^*.
\]
Thus
\[
\rho_m(5)=\frac1{4}
\]
for every \(m\in\F_5^*\), and
\[
\Delta_m(5)=0.
\]
The moment identity gives
\[
\frac1{|\Gamma_5|}-\frac1{5-1}
=
\frac14-\frac14
=
0.
\]

\subsection{A proper subgroup obstruction}

For \(q=23\), direct computation gives
\[
|\langle2,3\rangle\bmod23|=11.
\]
Thus the local density is
\[
\rho_m(23)=
\begin{cases}
1/11,& -m^{-1}\in\langle2,3\rangle,\\
0,& -m^{-1}\notin\langle2,3\rangle.
\end{cases}
\]
The second moment is exactly
\[
\frac1{11}-\frac1{22}
=
\frac1{22}.
\]
This example shows the obstruction energy directly: the subgroup has index \(2\), and the centered variance records that index.

\section{Further directions}

The calculus above is elementary, but it supplies a reusable foundation for density-zero and conditional Kummer-distribution results in problems involving \(m a^k b^\ell+1\).  The analytic problem of proving strong enough prime distribution in Kummer families is deliberately not claimed here.

\begin{thebibliography}{9}

\bibitem{Apostol}
T. M. Apostol.
\newblock \emph{Introduction to Analytic Number Theory}.
\newblock Springer, 1976.

\bibitem{DF}
D. S. Dummit and R. M. Foote.
\newblock \emph{Abstract Algebra}.
\newblock Wiley, 3rd edition, 2004.

\bibitem{IK}
H. Iwaniec and E. Kowalski.
\newblock \emph{Analytic Number Theory}.
\newblock American Mathematical Society Colloquium Publications, 2004.

\end{thebibliography}

\end{document}
